\section*{Chapter \ref{ch:b3:asymmetric-synthesis} --- Asymmetric Synthesis}

\begin{solution}{exo:b3:asymmetric-synthesis:1}
ee 80\,\%: $\mathrm{er} = 1.8/0.2 = 9$ (90\,:\,10). er 99\,:\,1: $\mathrm{ee} = 98/100 = 98$\,\%. ee 50\,\%: 75\,\% and 25\,\%.
\end{solution}

\begin{solution}{exo:b3:asymmetric-synthesis:2}
$(1840 - 60)/(1840 + 60) = 0.937$: ee 93.7\,\%, an er of $1840/60 = 31$.
\end{solution}

\begin{solution}{exo:b3:asymmetric-synthesis:3}
Priorities O $>$ ethyl $>$ H. Drawn with O up, ethyl lower left and H lower right, the
front face is \textit{Si}, the back face \textit{Re}. A methyl group added from the \omterm{def:b3:asymmetric-synthesis:faces}{Re face}, with
priorities OH $>$ ethyl $>$ methyl $>$ H in the product, gives (\textit{R})-butan-2-ol.
\end{solution}

\begin{solution}{exo:b3:asymmetric-synthesis:4}
Ethanol: \omterm{def:b3:asymmetric-synthesis:topicity}{enantiotopic} (exchanged by the \omterm{def:b3:point-groups:mirror}{mirror plane} of the molecule; same
shift in an achiral solvent). Butan-2-ol: \omterm{def:b3:asymmetric-synthesis:topicity}{diastereotopic} (next to a stereocentre;
different shifts).
\end{solution}

\begin{solution}{exo:b3:asymmetric-synthesis:5}
er 199: $RT\ln 199 = \qty{13.1}{kJ/mol}$ at \qty{298}{K}, \qty{8.6}{kJ/mol} at \qty{195}{K}. 80\,\% ee at \qty{25}{\celsius} is er 9,
$\Delta\Delta G^\ddagger = RT\ln 9 = \qty{5.45}{kJ/mol}$; at \qty{-78}{\celsius}, $\mathrm{er} = \eu^{5450/(8.314 \times 195.15)} = 28.7$, ee 93\,\%.
\end{solution}

\begin{solution}{exo:b3:asymmetric-synthesis:6}
\omterm{def:b3:asymmetric-synthesis:ee}{Optical purity} $12.0/15.0 = 80$\,\%: 90\,\% (\textit{S}) and 10\,\% (\textit{R}). A small rotation, impurities,
concentration and solvent effects, and non-linear rotation--composition
behaviour make it unreliable; chromatography separates and counts the
enantiomers directly.
\end{solution}

\begin{solution}{exo:b3:asymmetric-synthesis:7}
$s = \ln(0.40 \times 0.10)/\ln(0.40 \times 1.90) = -3.22/-0.274 = 12$.
\end{solution}

\begin{solution}{exo:b3:asymmetric-synthesis:8}
With L-($+$)-diethyl tartrate, (2\textit{S},3\textit{S})-2,3-epoxyhexan-1-ol; with D-($-$)-diethyl tartrate, the
enantiomer, (2\textit{R},3\textit{R}).
\end{solution}

\begin{solution}{exo:b3:asymmetric-synthesis:9}
L = phenyl, M = methyl, S = H. For the (\textit{S}) aldehyde the reactive conformer has phenyl
perpendicular to the carbonyl plane and the hydrogen beside the path of the
nucleophile, which attacks anti to phenyl. Working out the configurations gives
(2\textit{S},3\textit{S})-3-phenylbutan-2-ol as the major diastereomer.
\end{solution}

\begin{solution}{exo:b3:asymmetric-synthesis:10}
At the start the two enantiomers are present in equal amounts, so the products
form in the ratio $k_{\mathrm{fast}} : k_{\mathrm{slow}} = s$: $\mathrm{ee}_p = (s - 1)/(s + 1)$, 0.905 for $s = 20$. As the reaction goes on the
fast enantiomer is depleted and the product ee falls, while the substrate ee rises
towards 1: the substrate can always be purified by going further; the product
cannot.
\end{solution}

\begin{solution}{exo:b3:asymmetric-synthesis:11}
With fast racemisation: up to 100\,\% yield, $\mathrm{er} = 99$, ee 98\,\%. Without it, at 50\,\%
conversion: yield at most 50\,\%, and a product ee of 93\,\% (the substrate's ee is the
same, 93\,\%).
\end{solution}

\begin{solution}{exo:b3:asymmetric-synthesis:12}
When the two complexes interconvert faster than they react, the product ratio is
$k_B[\mathrm B]/(k_A[\mathrm A]) = 600/10 = 60$: the minor complex B gives 98\,\% of the product (ee 97\,\%). The
selectivity is set by the difference between the two transition states, not by
the \omterm{def:b3:partition-functions:boltzmann}{populations} of the intermediates.
\end{solution}

\begin{solution}{pb:b3:asymmetric-synthesis:1}
\textbf{1.} The alkene carbon bearing the nitrogen has three different substituents
and two different faces; hydrogen added to it makes the $\alpha$-carbon of the amino
acid a stereocentre.

\textbf{2.} The racemate: both faces are attacked at the same rate.

\textbf{3.} Its carbonyl oxygen binds the rhodium with the C=C bond: the substrate is
held as a chelate, in a fixed geometry that the ligand can discriminate.

\textbf{4.} It builds a chiral pocket around the metal ($C_2$-symmetric): the two ways of
binding the substrate, by one face or the other, are diastereomeric, with different
energies and reactivities.

\textbf{5.} No: the minor complex reacts much faster with hydrogen (Curtin--Hammett).

\textbf{6.} The metal is on one face of the bound alkene, and the hydrogens are
transferred from the metal.

\textbf{7.} $\mathrm{er} = 1.95/0.05 = 39$.

\textbf{8.} $\Delta\Delta G^\ddagger = RT\ln 39 = 8.314 \times 298.15 \times 3.664 = \qty{9.1}{kJ/mol}$.

\textbf{9.} $\mathrm{er} = \eu^{9082/(8.314 \times 195.15)} = 270$: ee 99.3\,\%.

\textbf{10.} $\mathrm{er} = \eu^{11082/(8.314 \times 298.15)} = 87$: ee 97.7\,\%.

\textbf{11.} Less than a typical hydrogen bond: a small difference in the steric
contacts or one weak interaction between substrate and ligand decides the
outcome, which is why such catalysts are hard to design from first principles.

\textbf{12.} The rate falls, hydrogen dissolves and the catalyst behaves differently
at low temperature, and cooling a large reactor costs energy.

\textbf{13.} 97.5 g of the (\textit{S}) and 2.5 g of the (\textit{R}) enantiomer.

\textbf{14.} $97.5 - 5.0 = \qty{92.5}{g}$ of crystals, practically enantiopure.

\textbf{15.} $92.5/100 = 92.5$\,\% of the crude (95\,\% of the (\textit{S}) enantiomer).

\textbf{16.} The crystals take the excess enantiomer, the racemic part staying in
solution with the minor one; from a racemate there is no excess to collect, and
crystallisation gives racemic crystals.

\textbf{17.} By chiral HPLC against a racemic reference.

\textbf{18.} 50\,\%: only the slow enantiomer survives.

\textbf{19.} $s = \ln(0.45 \times 0.05)/\ln(0.45 \times 1.95) = -3.79/-0.131 = 29$.

\textbf{20.} 45\,\% of the racemate (88\,\% of the desired enantiomer present at the start,
at 95\,\% ee).

\textbf{21.} With fast racemisation of the substrate the yield could approach 100\,\%.

\textbf{22.} It converts all of the cheap achiral precursor into the right enantiomer,
with a catalyst used in tiny amounts and no enantiomer thrown away.

\textbf{23.} Rhodium is toxic and its residues in drugs are limited to a few parts per
million; they are measured by inductively coupled plasma mass spectrometry.

\textbf{24.} 95\,\% ee at \qty{25}{\celsius} corresponds to $\Delta\Delta G^\ddagger = RT\ln 39 \approx \qty{9.1}{kJ/mol}$.
\end{solution}
